\setcounter{section}{7}

\Needspace{5\baselineskip}
\hypertarget{ux5e8fux5217ux5230ux5e8fux5217seq2sequx5b66ux4e60}{%
\section{序列到序列（Seq2Seq）学习}\label{ux5e8fux5217ux5230ux5e8fux5217seq2sequx5b66ux4e60}}

从技术上讲，编码器将长度可变的输入序列转换成形状固定的上下文变量 \(\mathbf{c}\)，并且将输入序列的信息在该上下文变量中进行编码。可以使用循环神经网络来设计编码器。

\Needspace{5\baselineskip}
考虑由一个序列组成的样本（批量大小是 1）。假设输入序列是 \(x_1,\ldots,x_T\)，其中 \(x_t\) 是输入文本序列中的第 \(t\) 个词元。在时间步 \(t\)，循环神经网络将词元 \(x_t\) 的输入特征向量 \(\mathbf{x}_t\) 和 \(\mathbf{h}_{t-1}\)（即上一时间步的隐状态）转换为 \(\mathbf{h}_t\)（即当前步的隐状态）。使用一个函数 \(f\) 来描述循环神经网络的循环层所做的变换：

\[
\mathbf{h}_t = f(\mathbf{x}_t,\mathbf{h}_{t-1})
\]

\Needspace{5\baselineskip}
总之，编码器通过选定的函数 \(q\)，将所有时间步的隐状态转换为上下文变量：

\[
\mathbf{c} = q(\mathbf{h}_1,\ldots,\mathbf{h}_T)
\]

正如上文提到的，编码器输出的上下文变量 \(\mathbf{c}\) 对整个输入序列 \(x_1,\ldots,x_T\) 进行编码。来自训练数据集的输出序列 \(y_1,y_2,\ldots,y_{T'}\)，对于每个时间步 \(t'\)（与输入序列或编码器的时间步 \(t\) 不同），解码器输出 \(y_{t'}\) 的概率取决于先前的输出子序列 \(y_1,\ldots,y_{t'-1}\) 和上下文变量 \(\mathbf{c}\)，即 \(P(y_{t'} \mid y_1,\ldots,y_{t'-1},\mathbf{c})\)。

\Needspace{5\baselineskip}
为了在序列上模型化这种条件概率，我们可以使用另一个循环神经网络作为解码器。在输出序列上的任意时间步 \(t'\)，循环神经网络将来自上一时间步的输出 \(y_{t'-1}\) 和上下文变量 \(\mathbf{c}\) 作为其输入，然后在当前时间步将它们和上一隐状态 \(\mathbf{s}_{t'-1}\) 转换为隐状态 \(\mathbf{s}_{t'}\)。因此，可以使用函数 \(g\) 表示解码器的隐藏层的变换：

\[
\mathbf{s}_{t'} = g(y_{t'-1},\mathbf{c},\mathbf{s}_{t'-1})
\]

在获得解码器的隐状态之后，我们可以使用输出层和 softmax 操作来计算在时间步 \(t'\) 时输出 \(y_{t'}\) 的条件概率分布 \(P(y_{t'} \mid y_1,\ldots,y_{t'-1},\mathbf{c})\)。

在每个时间步，解码器预测了输出词元的概率分布。类似于语言模型，可以使用softmax来获得分布，并通过计算交叉熵损失函数来进行优化。另外，为了让不同长度的序列可以以相同形状的小批量加载。 特定的填充词元被添加到序列的末尾， 我们应该将填充词元的预测排除在损失函数的计算之外。用sequence\_mask函数可通过零值化屏蔽不相关的项， 以便后面任何不相关预测的计算都是与零的乘积，结果都等于零（交叉熵损失-sigma(pi*logqi),pi=0）。

\FloatBarrier
\Needspace{5\baselineskip}
\hypertarget{ux53c2ux8003ux6587ux732e-24}{%
\subsection{参考文献}\label{ux53c2ux8003ux6587ux732e-24}}

\vspace{0.25em}
\begin{itemize}
\tightlist
\item
  Zhang, A., Lipton, Z. C., Li, M., \& Smola, A. J. (2023). \href{https://D2L.ai}{Dive into Deep Learning}. Cambridge University Press.
\item
  Sutskever, I., Vinyals, O., \& Le, Q. V. (2014). \href{https://proceedings.neurips.cc/paper_files/paper/2014/hash/5a18e133cbf9f257297f410bb7eca942-Abstract.html}{Sequence to Sequence Learning with Neural Networks}. NeurIPS 2014.
\end{itemize}

\clearpage
